% ============================================================
%  ch03a  §3 引言 + §3.1 志村簇的局部模型
%  原文对应：ICM-arXiv.tex 1070--1105 行
%  编号：公式 (3.1)；定理 3.1.1；注 3.1.2
% ============================================================

\section{对算术几何的应用}\label{Application}

除前面提到的几何朗兰兹（的想法）对经典朗兰兹纲领的直接应用之外，我们再讨论一些算术应用，它们大多与志村簇有关，且基于作者本人的工作。需要指出，几何朗兰兹（的想法）对算术问题还有许多其他值得注意的应用，例如 \cite{Gaitsgory-deJong, HNY-Kloosterman, Yun-motive-exceptional-group, Xu-Zhu-Bessel, LLLM}，这里仅举几例。

\subsection{志村簇的局部模型}\label{S: local model}

志村簇整模型（带抛孔子群水平）的理论，始于（隐含在 Kronecker 工作中的）对带 $\Ga_0(p)$-水平的椭圆模曲线的 mod $p$ 约化的理解。我们讨论这一理论的一小部分，即借助局部模型理论研究这些整模型的 étale 局部结构。局部模型理论的近期发展深受几何朗兰兹纲领的影响。

沿用 \S\ref{geometric JL} 中关于志村簇的记号（但本小节不假设 $G$ 为伴随型）。设 $(G,X)$ 为志村数据，$K$ 为所选水平，满足 $K_p=\mG(\bZ_p)$，其中 $\mG$ 是 $G_{\bQ_p}$ 在 $\bZ_p$ 上的某个抛孔子群概形（Bruhat--Tits 意义下）。于是对 $E$ 上位于 $p$ 之上的位置 $v$，局部模型图指 $\mO_{E_v}$ 上拟射影概形之间的对应
\begin{equation}\label{E: local model diagram}
\scrS_K\leftarrow\widetilde{\scrS}_K\xrightarrow{\tilde\varphi} M^{\loc}_{\mG},
\end{equation}
其中 $\scrS_K$ 是 $\Sh_K(G)$ 在 $\mO_{E_v}$ 上的整模型，$\widetilde{\scrS}_K$ 是 $\scrS_K$ 上的 $\mG_{\mO_{E_v}}$-挠子，$M^{\loc}_{\mG}$ 是所谓的局部模型，即 $\mO_{E_v}$ 上带 $\mG_{\mO_{E_v}}$-作用的平坦射影概形，而 $\tilde{\varphi}$ 是相对维数为 $\dim G$ 的 $\mG_{\mO_{E_v}}$-等变光滑态射。因此 $M_{\mG}^{\loc}$ 刻画了 $\scrS_K$ 的 étale 局部结构。另一方面，$M_{\mG}^\loc$ 上 $\mG_{\mO_{E_v}}$-作用的存在，使它比 $\scrS_K$ 更易研究。

局部模型最初的构造，基于把抛孔子群概形实现为（$F$ 上除代数中）带双线性型的向量空间里一条自对偶格链的稳定化群（的幺单位连通分支），综述与文献见 \cite{PRS-Survey}。这一途径多少有些特设，且局限于所谓的（P）EL（局部）志村数据。基于 \S\ref{Tamely local} 中叠 $\Hk_{\mG,D}$ 的 $\bZ_p$-类比之构造的新途径，由 \cite{Pappas-Zhu-Local-Model} 系统引入（其中假设 $G$ 驯顺，该假设后来在 \cite{Levin-Weil-Restriction, Lourenco} 中被去掉）。在该文中，这种 $\bZ_p$-类比（更确切地说，是相应的 Beilinson--Drinfeld 型仿射 Grassmann 簇 $\Gr_{\mG,\bZ_p}$）的构造，基于对某个“二维抛孔子”群概形 $\widetilde\mG$ 的构造：它是 $\bZ_p[\varpi]$ 上的群概形，沿 $\bZ_p[\varpi]\xrightarrow{\varpi\mapsto p}\bZ_p$ 的限制可还原出 $\mG$。（综述见 \cite{Zhu-ICCM}。）

这种仿射 Grassmann 簇的另一 $p$-进版本 $\Gr_{\mG,\Spd\bZ_p}$，在 \cite{Scholze-Berkeley} 中于解析的完满世界中给出了更直接的构造。无论哪种情形，局部模型都定义为生成纤维中对应于 $\mu$ 的 Schubert 簇的平坦闭包。此外，近期工作 \cite{AGLR} 表明这两种构造一致。

下述定理取自 \cite{AGLR}，是关于局部模型存在性及其性质的最新结果。

\begin{thm}
设 $G$ 是 $p$-进域 $F$ 上的连通约化群。除 $p=2$ 的奇酉情形与 $p=3$ 的三对偶情形之外，对 $G$ 在 $\mO$ 上的任一抛孔子群概形 $\mG$，以及 $G$ 的、定义在 $F$ 的有限扩张 $E/F$ 上的微小余权共轭类 $\mu$，都存在 $\mO_E$ 上的正规平坦射影概形 $M_{\mG,\mu}^{\loc}$，带有 $\mG_{\mO_E}$-作用，使得 $M_{\mG,\mu}^{\loc}\otimes E$ 与 $G_E$ 中对应于 $\mu$ 的部分旗簇 $\Fl_{\mu}$ 是 $G_E$-等变同构的，而 $M_{\mG}^{\loc}\otimes k_E$ 与 $LG/L^+\mG\otimes k_E$ 中 $\mu$-容许集上的 Schubert 簇之并（的典范去完美化）是 $(\mG\otimes k_E)$-等变同构的。此外，$M_{\mG}^{\loc}$ 正规、Cohen--Macaulay，且其特殊纤维的每个几何不可约分支都正规且 Cohen--Macaulay。
\end{thm}

本小节以几条注记结束。

\begin{rmk}
\begin{enumerate}
\item 一旦建立局部模型图 \eqref{E: local model diagram}，该定理也就给出志村簇整模型的相应性质。

\item 研究局部模型特殊纤维的一个关键要素，是 Pappas--Rapoport 的凝聚猜想 \cite{Pappas-Rapoport-Ramified-Unitary}，它由 \cite{Zhu-Coherence-Conjecture} 证明（证明用到了融合积的想法）。

\item 局部模型理论的一个重要动机／应用是 Haines--Kottwitz 猜想 \cite{Haines-Test-Function}：它预测，抛孔 Hecke 代数 $H_{K_p}^{\cl}$ 中的某个中心元应当作为计算 $\Sh_{K}(G)$ 的 Hasse--Weil zeta 函数的迹公式中的测试函数。如 \S\ref{Tamely local} 所述，这一猜想正是 Gaitsgory 中心层构造 \eqref{E: central sheaf} 的动机之一。在构造出 $\bZ_p$ 上的局部 Hecke 叠 $\Hk_{\mG,\bZ_p}$（无论用 \cite{Pappas-Zhu-Local-Model} 还是 \cite{Scholze-Berkeley} 的版本）之后，就可以在混合特征下模仿构造 \eqref{E: central sheaf}，从而解决 Kottwitz 猜想。最新结果同样见 \cite{AGLR}。
\end{enumerate}
\end{rmk}
